\documentclass[11pt]{article}
\usepackage[margin=1in]{geometry}
\usepackage{amsmath,amssymb}
\usepackage{physics}
\usepackage{hyperref}

\newcommand{\Gam}{\Gamma}
\newcommand{\eps}{\varepsilon}
\newcommand{\BZ}{\mathrm{BZ}}
\newcommand{\MBZ}{\mathrm{MBZ}}
\newcommand{\Nk}{N_k}
\newcommand{\Ang}{\text{\AA}}

\title{Self-Consistent Born Approximation (SCBA) \\
for Hofstadter Transport Coefficients}
\author{Notes for internal reference}
\date{}

\begin{document}
\maketitle

%======================================================================
\section{Formalism}
%======================================================================

\subsection{Kubo formula for longitudinal conductivity}

The dissipative (longitudinal) conductivity at chemical potential $\mu$
and temperature $T$ is given by the Kubo--Bastin formula in the
spectral-function representation~\cite{Bastin1971,Streda1982}:
\begin{equation}\label{eq:sxx_spectral}
  \sigma_{xx}(\mu) = \frac{\pi e^2 \hbar}{\Omega}
  \int d\eps\, \Bigl(-\frac{\partial f}{\partial\eps}\Bigr)
  \Phi_{xx}(\eps),
\end{equation}
where $\Omega$ is the system area, $f(\eps)$ is the Fermi--Dirac
distribution at temperature $T$, and the spectral conductivity is
\begin{equation}\label{eq:Phi_xx}
  \Phi_{xx}(\eps) = \sum_{\mathbf{k},n,m}
  |v^x_{nm}(\mathbf{k})|^2\, A_n(\eps,\mathbf{k})\, A_m(\eps,\mathbf{k}).
\end{equation}
Here $v^x_{nm} = \langle n\mathbf{k}|\hat{v}_x|m\mathbf{k}\rangle$
are velocity matrix elements in the eigenbasis of $H(\mathbf{k})$,
and the spectral function of eigenstate $n$ is
\begin{equation}\label{eq:spectral_fn}
  A_n(\eps) = \frac{1}{\pi}\,
  \frac{\Gam(\eps)}{(\eps - E_n)^2 + \Gam(\eps)^2}.
\end{equation}
For constant broadening, $\Gam(\eps) = \Gam_0$ is
energy-independent.  Substituting Eq.~\eqref{eq:spectral_fn} into
\eqref{eq:Phi_xx} and pulling out the $\Gam^2/\pi^2$ prefactor:
\begin{equation}\label{eq:Phi_xx_L}
  \Phi_{xx}(\eps) = \frac{\Gam(\eps)^2}{\pi^2}
  \sum_{\mathbf{k},n,m} |v^x_{nm}|^2\, L_n(\eps)\, L_m(\eps),
\end{equation}
where $L_n(\eps) \equiv [(\eps - E_n)^2 + \Gam(\eps)^2]^{-1}$ is the
(unnormalized) Lorentzian.  At $T=0$,
$-\partial f/\partial\eps = \delta(\eps - \mu)$ and
\begin{equation}
  \sigma_{xx}(\mu)\Big|_{T=0}
  = \frac{e^2 \hbar}{\pi\,\Omega}\, \Gam(\mu)^2
  \sum_{\mathbf{k},n,m} |v^x_{nm}|^2\, L_n(\mu)\, L_m(\mu).
\end{equation}

\subsection{Hall conductivity}

The Hall conductivity is~\cite{TKNN1982,Streda1982}:
\begin{equation}\label{eq:sxy}
  \sigma_{xy}(\mu) = \frac{e^2}{\hbar}\,\frac{1}{2\pi}\,
  \frac{1}{\Omega}\sum_{\mathbf{k}} \sum_{E_n < \mu} \mathcal{K}_n(\mathbf{k}),
\end{equation}
where the broadened Berry curvature kernel is
\begin{equation}
  \mathcal{K}_n(\mathbf{k}) = \sum_{m \neq n}
  \frac{\mathrm{Im}[v^x_{nm}\, (v^y_{nm})^*]}
  {(E_n - E_m)^2 + \Gam(E_n)^2}.
\end{equation}
By default, $\Gam(E_n)$ is the SCBA broadening of the occupied
state~$n$.  This sharpens $\sigma_{xy}$ features at low flux where
subbands are narrower than~$\Gam_0$.  Chern quantization in spectral
gaps is preserved: $\Gam(E) \to \Gam_0 f_{\min}$ in gaps, and the
Chern number is insensitive to the magnitude of the regularizer.
Setting \texttt{scba\_xy\_constant = 1} reverts to constant~$\Gam_0$.
In the limit $\Gam \to 0$, $\sigma_{xy}$ reduces to the TKNN formula,
and its value in a spectral gap gives the Chern number~$C$.

\subsection{The SCBA self-consistency equation}

For short-range (white-noise) disorder with scattering strength
$\Gam_0$ (the bare broadening), the self-consistent Born
approximation~\cite{Ando1974} replaces the constant
$\Gam_0$ with an energy-dependent $\Gam(E)$ determined by:
\begin{equation}\label{eq:scba}
  \boxed{
  \Gam(E) = \pi\, \Gam_0^2\, \rho(E)
  }
\end{equation}
where the density of states per unit cell is computed with the
current $\Gam(E)$:
\begin{equation}\label{eq:rho}
  \rho(E) = \frac{1}{N_\BZ}\sum_{\mathbf{k},n}
  \frac{1}{\pi}\,\frac{\Gam(E)}{(E - E_{n\mathbf{k}})^2 + \Gam(E)^2}.
\end{equation}
Here $N_\BZ$ is the total number of unit cells in the BZ sum (so that
$\rho$ is a DOS \emph{per unit cell}), and the sum runs over all
bands~$n$ and $\mathbf{k}$-points.

\textbf{Physical content.}  Equation~\eqref{eq:scba} makes $\Gam(E)$
proportional to the local DOS.  In spectral gaps, $\rho \to 0$ so
$\Gam \to 0$: states are not broadened into the gap.  In narrow
Hofstadter subbands with bandwidth $W \ll \Gam_0$, the self-consistent
$\Gam(E)$ shrinks below $\Gam_0$, reflecting Anderson localization at
mean-field level.  The dissipative conductivity
$\sigma_{xx} \propto \Gam^2\, \rho^2$ is then strongly suppressed in
narrow bands.

\textbf{SCBA broadening in $\sigma_{xy}$.}  By default the code uses
$\Gam(E_n)$ in the Berry curvature kernel.  Since $\Gam(E)$ is small
(floored at $f_{\min}\,\Gam_0$) in spectral gaps, the Chern
quantization $\sigma_{xy} = (e^2/h)\,C$ is preserved---indeed
sharpened relative to constant $\Gam_0$.  At low flux, where Hofstadter
subbands are narrower than $\Gam_0$, this produces $\sigma_{xy}$
features that are consistent with the SCBA-sharpened $\sigma_{xx}$.

\subsection{Iterative solution}

Equations \eqref{eq:scba}--\eqref{eq:rho} are solved iteratively:
\begin{enumerate}
  \item Initialize $\Gam^{(0)}(E) = \Gam_0$ on a fine energy grid.
  \item Compute $\rho^{(i)}(E)$ from the current $\Gam^{(i)}(E)$
    via Eq.~\eqref{eq:rho}.
  \item Compute the ``target'' $\widetilde{\Gam}^{(i)}(E)
    = \pi\,\Gam_0^2\, \rho^{(i)}(E)$.
  \item Enforce a floor:
    $\widetilde{\Gam}^{(i)}(E) \ge f_{\min}\,\Gam_0$ to prevent
    numerical instability in spectral gaps.
  \item Update $\Gam^{(i+1)}$ via Anderson (Pulay) mixing of
    depth~$M$, or linear mixing as a fallback.
  \item Converge when
    $\max|\widetilde{\Gam} - \Gam^{(i)}| / \max(\Gam^{(i)}) < \tau$.
\end{enumerate}

\paragraph{Anderson mixing.}
Given the last $M$ iterates $\{x_j\}$ and residuals
$\{R_j = g(x_j) - x_j\}$, form the matrix
$\Delta R = [R_{i} - R_j]$ and solve
$\Delta R\, \mathbf{c} = R_{\mathrm{current}}$ in the least-squares
sense.  The new iterate is
$x_{\mathrm{new}} = \widetilde{x} - \Delta X\,\mathbf{c}$, where
$\Delta X = [\widetilde{x} - x_j]$.  This is equivalent to a
multi-secant (quasi-Newton) update and typically converges in
$\sim$25--30 iterations versus $\sim$200 for linear mixing.

%======================================================================
\section{Implementation in the Hofstadter code}
%======================================================================

This section describes how the formalism of Section~1 maps onto the
code in \texttt{main\_v3.py}.

\subsection{Brillouin zone geometry and normalization}

The moire superlattice has a ``primary'' unit cell area
\begin{equation}
  A_{\mathrm{prim}} = \frac{\sqrt{3}}{2}\, L_{\mathrm{moire}}^2,
\end{equation}
and the magnetic flux per primitive cell is
$\Phi/\Phi_0 = qq/(2\,pp)$, i.e.
$B = \bigl(qq/(2\,pp)\bigr)\,\Phi_0 / A_{\mathrm{prim}}$.
(The input fraction $qq/pp$ is the flux through the
centered-rectangular two-lattice-point cell of area
$2\,A_{\mathrm{prim}}$ used internally by the Landau-gauge
construction.)

The magnetic unit cell --- the region whose states appear exactly once
in the spectrum at each $\mathbf{k}$-point --- contains $2\,pp$
primitive cells (area $A_{\mathrm{mag}} = 2\,pp\,A_{\mathrm{prim}}$),
enclosing $qq$ flux quanta, matching the $qq$ guiding centers per
Landau level in the basis.

Eigenvalues are computed on a uniform $\mathbf{k}$-mesh of $\Nk$
points.  The \emph{physical} DOS per primitive unit cell requires
dividing the eigenvalue sum by the number of primitive unit cells
represented:
\begin{equation}\label{eq:Nbz}
  N_\BZ = \Nk \times 2\,pp.
\end{equation}
The factor $2\,pp$ is the number of primitive cells per magnetic unit
cell (one magnetic cell of states per $\mathbf{k}$-point).

\subsection{SCBA solver: \texttt{\_solve\_scba()}}

\paragraph{Inputs.}  The function receives:
\begin{itemize}
  \item \texttt{all\_eigs\_eV}: array of shape
    $(\Nk \times n_{\mathrm{valleys}},\; n_b)$ containing all
    eigenvalues from both valleys, in~eV.  Both K and K$'$ eigenvalues
    are included because white-noise disorder scatters between valleys.
  \item \texttt{Gamma0}: the bare broadening $\Gam_0$ in eV.
  \item \texttt{pp}, \texttt{Nk}: for the normalization factor.
\end{itemize}

\paragraph{Energy grid.}  A uniform grid of spacing
$\delta E = \Gam_0/10$ is constructed over the eigenvalue range plus a
margin of $5\Gam_0 + 10\%$ of the bandwidth.

\paragraph{DOS computation.}  At each iteration, the DOS is computed as
\begin{equation}
  \rho(E_i) = \frac{1}{\pi\,\Nk\, 2\,pp}
  \sum_{j} \frac{\Gam(E_i)}{(E_i - E_j)^2 + \Gam(E_i)^2},
\end{equation}
where $j$ runs over all eigenvalues in the flattened
\texttt{all\_eigs\_eV} array.  This is line~221 of \texttt{main\_v3.py}:
\begin{verbatim}
  dos_norm = 1.0 / (np.pi * Nk * 2 * pp)
\end{verbatim}
The computation is chunked over energy grid points to limit memory
usage (at most $\sim$50M floats per chunk).

\paragraph{Update and mixing.}  The target broadening
$\widetilde{\Gam} = \pi\,\Gam_0^2\,\rho$ is computed (line~238),
floored (line~239), and then mixed via Anderson acceleration of
depth~$M$ (default 5), falling back to linear mixing with
$\alpha = 0.3$ on the first iteration or when the least-squares system
is singular.

\subsection{Two-pass k-loop}

Because the SCBA needs all eigenvalues to solve for $\Gam(E)$, and
the transport Kubo formulas need $\Gam(E)$ already known, the code
uses two passes:

\begin{enumerate}
  \item \textbf{Pass 1 (eigenvalue collection):}
    \texttt{\_solve\_kpoint\_eigenvalues\_core()} diagonalizes $H(\mathbf{k})$
    using \texttt{eigvalsh} (eigenvalues only, no eigenvectors).
    The eigenvalues from all $\mathbf{k}$-points and both valleys are
    collected into \texttt{all\_eigs\_eV} and passed to
    \texttt{\_solve\_scba()}.

  \item \textbf{Pass 2 (transport):}
    \texttt{\_solve\_kpoint\_transport\_core()} diagonalizes
    $H(\mathbf{k})$ again with \texttt{eigh} (eigenvalues \emph{and}
    eigenvectors), computes velocity matrix elements, and evaluates the
    Kubo sums using the now-known $\Gam(E)$.
\end{enumerate}

\subsection{Effect on $\sigma_{xx}$}

\paragraph{Spectral function.}  In SCBA mode, the spectral function
uses $\Gam(\eps)$ evaluated at the \emph{probe} energy~$\eps$:
\begin{equation}
  A_n(\eps) = \frac{1}{\pi}\,
  \frac{\Gam(\eps)}{(\eps - E_n)^2 + \Gam(\eps)^2}.
\end{equation}
Because $\Gam$ depends on $\eps$ (not on eigenstate indices $n,m$),
the $\Gam^2/\pi^2$ factor in Eq.~\eqref{eq:Phi_xx_L} can still be
pulled out of the $n,m$ sum at each grid point:
\begin{equation}
  \Phi_{xx}(\eps) = \frac{\Gam(\eps)^2}{\pi^2}
  \sum_{n,m} |v^x_{nm}|^2\, L_n(\eps)\, L_m(\eps),
\end{equation}
with $L_n(\eps) = [(\eps - E_n)^2 + \Gam(\eps)^2]^{-1}$.

\paragraph{Code structure.}  In \texttt{\_transport\_kubo\_single\_k()},
the SCBA branch (lines~151--156) interpolates $\Gam(\eps)$ from the
precomputed grid, builds the Lorentzians $L_n(\eps)$ with
$\Gam(\eps)$, and multiplies the result by $\Gam(\eps)^2$:
\begin{verbatim}
  G_eps = np.interp(eps_grid, scba_grid, scba_Gamma)
  G2_eps = G_eps ** 2
  L_all = 1.0 / ((eps_grid[:,None] - E[None,:])**2 + G2_eps[:,None])
  Phi_xx = np.sum((L_all @ vx_sq) * L_all, axis=1) * G2_eps
\end{verbatim}

\paragraph{Prefactor.}  The caller applies:
\begin{equation}
  \sigma_{xx} = \frac{2\,\hbar^2}{A_{\mathrm{mag}}\,\Nk}
  \times \sum_\mathbf{k} \Phi_{xx}(\mu) \quad
  [\text{units of } e^2/h],
\end{equation}
where $A_{\mathrm{mag}} = 2\,pp\,A_{\mathrm{prim}}$ is the magnetic
unit cell area in~$\Ang^2$ (one magnetic cell of states per
$\mathbf{k}$-point; no $qq$ appears in the normalization).  The
output is per spin and per valley, matching the $\sigma_{xy}$
convention; the factor of 2 comes from the $\pi$ factors in the
spectral functions (see derivation below).  In SCBA mode the $\Gam^2$
is already inside $\Phi_{xx}$, so \texttt{pf\_xx} omits it:
\begin{verbatim}
  pf_xx = 2.0 * HBAR_EV**2 / (A_mag_Ang2 * Nk_tot)   # SCBA
  pf_xx = 2.0 * HBAR_EV**2 * G2 / (A_mag_Ang2 * Nk_tot)  # constant
\end{verbatim}

\paragraph{Prefactor derivation.}  Starting from Eq.~\eqref{eq:sxx_spectral}
at $T=0$ (one spin, one valley, one $\mathbf{k}$-point):
\begin{align}
  \sigma_{xx} &= \frac{\pi e^2 \hbar}{\Omega}\, \Phi_{xx}(\mu)
  = \frac{\pi e^2 \hbar}{\Omega}\, \frac{\Gam^2}{\pi^2}\,
  \sum_{nm} |v^x_{nm}|^2\, L_n\, L_m \nonumber \\
  &= \frac{e^2 \hbar}{\pi\,\Omega}\, \Gam^2\,
  \sum_{nm} |v^x_{nm}|^2\, L_n\, L_m.
\end{align}
To express in units of $e^2/h$, divide by $e^2/h = e^2/(2\pi\hbar)$:
\begin{equation}
  \frac{\sigma_{xx}}{e^2/h}
  = \frac{2\hbar^2}{\Omega}\, \Gam^2\, \sum_{nm} |v^x_{nm}|^2\, L_n\, L_m.
\end{equation}
Including the k-mesh average
($\Omega \to A_{\mathrm{mag}}$, the area whose states appear once per
$\mathbf{k}$-point, with the $1/\Nk$ sum over the mesh), and keeping
the result per spin and per valley (no spin degeneracy factor):
\begin{equation}
  \frac{\sigma_{xx}}{e^2/h}
  = \frac{2\, \hbar^2}{A_{\mathrm{mag}}\,\Nk}\, \Gam^2\,
  \sum_{\mathbf{k},n,m} |v^x_{nm}|^2\, L_n\, L_m.
\end{equation}
In SCBA mode, $\Gam^2$ is absorbed into the per-$\eps$ sum (it varies
with $\eps$), leaving the prefactor without $\Gam^2$.

\subsection{Effect on $\sigma_{xy}$}

By default, the Berry curvature kernel uses $\Gam(E_n)$:
\begin{equation}
  \mathcal{K}_n = \sum_{m \neq n}
  \frac{\mathrm{Im}[v^x_{nm}\,(v^y_{nm})^*]}{(E_n - E_m)^2 + \Gam(E_n)^2}.
\end{equation}
The broadening $\Gam(E_n)$ of the occupied state~$n$ regularizes the
interband denominator.  In spectral gaps, $\Gam(E_n)$ drops to the
floor $f_{\min}\,\Gam_0$, sharpening features without affecting
Chern quantization.  Setting \texttt{scba\_xy\_constant = 1} reverts
to constant~$\Gam_0$ in the denominator.

\subsection{Units summary}

\begin{center}
\begin{tabular}{lll}
  \hline
  Quantity & Internal units & Notes \\
  \hline
  Eigenvalues (k-loop) & meV & converted to eV for Kubo \\
  $\Gam_0$, $\Gam(E)$ & eV (internal), meV (output) & \\
  Velocity $v_{nm}$ & $\Ang$/s & $A$ in $\Ang$, $H$ in eV, $\hbar$ in eV$\cdot$s \\
  $A_m$ (mag.\ cell area) & $\Ang^2$ & $= pp^2\, A_{\mathrm{uc}} / (2\,qq) \times 10^{20}$ \\
  $\hbar$ for transport & eV$\cdot$s & \texttt{HBAR\_EV} $= 6.582 \times 10^{-16}$ \\
  $\sigma_{xx}$, $\sigma_{xy}$ & $e^2/h$ & dimensionless \\
  $L_{12,xx}$, $L_{12,xy}$ & $e^2/h \cdot \mathrm{eV}$ & \\
  \hline
\end{tabular}
\end{center}

\subsection{Input parameters}

\begin{center}
\begin{tabular}{llp{7cm}}
  \hline
  Parameter & Default & Role \\
  \hline
  \texttt{broadening} & \texttt{'constant'} & Set to \texttt{'scba'} to enable \\
  \texttt{Gamma} & 1.0\,meV & Bare broadening $\Gam_0$ \\
  \texttt{scba\_mixing} & 0.3 & Linear mixing parameter $\alpha$ \\
  \texttt{scba\_tol} & $10^{-4}$ & Convergence tolerance $\tau$ \\
  \texttt{scba\_maxiter} & 200 & Maximum iterations \\
  \texttt{scba\_floor} & 0.01 & Floor ratio $f_{\min}$ ($\Gam \ge f_{\min}\,\Gam_0$) \\
  \texttt{scba\_anderson} & 5 & Anderson mixing depth $M$ (0 = linear only) \\
  \hline
\end{tabular}
\end{center}

%======================================================================

\begin{thebibliography}{9}

\bibitem{Ando1974}
  T.~Ando,
  ``Theory of Quantum Transport in a Two-Dimensional Electron System
  under Magnetic Fields. II. Single-Site Approximation (SCBA),''
  J.~Phys.\ Soc.\ Jpn.\ \textbf{36}, 1521 (1974).

\bibitem{MacDonald1984}
  A.~H.~MacDonald,
  ``Quantized Hall effect in a hexagonal periodic potential,''
  Phys.\ Rev.\ B \textbf{29}, 3057 (1984).

\bibitem{TKNN1982}
  D.~J.~Thouless, M.~Kohmoto, M.~P.~Nightingale, and M.~den~Nijs,
  ``Quantized Hall Conductance in a Two-Dimensional Periodic Potential,''
  Phys.\ Rev.\ Lett.\ \textbf{49}, 405 (1982).

\bibitem{Bastin1971}
  A.~Bastin, C.~Lewiner, O.~Betbeder-Matibet, and P.~Nozieres,
  ``Quantum oscillations of the Hall effect of a fermion gas
  with random impurity scattering,''
  J.~Phys.\ Chem.\ Solids \textbf{32}, 1811 (1971).

\bibitem{Streda1982}
  P.~St\v{r}eda,
  ``Theory of quantised Hall conductivity in two dimensions,''
  J.~Phys.\ C \textbf{15}, L717 (1982).

\end{thebibliography}

\end{document}
